
\input{0_uothesis_preamble}

\input{1_uothesis_prefatory_pages.tex}

\input{2_uothesis_begin_main_body_of_document.tex}

\input{introduction.tex}


\chapter{Evaluating evidence for co-geography in the \textit{Anopheles-Plasmodium} host-parasite system}

\bigskip
This chapter was published in the journal \textit{G3 | Genes | Genomes | Genetics} in 2024.
Andrew D. Kern and Peter L. Ralph are co-authors on this work. 
Co-authors and I conceptualized the study together and I performed all analyses; 
I wrote the manuscript with A.D.K and editorial assistance from P.L.R.
\bigskip
The citation for this publication is as folllows:

\fullcite{rehmann2024evaluating}
\bigskip

\input{co_geography.tex}

\newpage
\section{Bridge}

In Chapter II, I use continuous space simulation and machine learning to evaluate evidence for co-geography in
the \textit{Anopheles-Plasmodium} host-parasite system, ultimately finding evidence for common dispersal patterns between
the parasite and its mosquito vector. 
Identifying patterns of spatial population structure and dispersal in this system, however, led me to ask further 
questions about how spatial structure may be shaping the evolutionary dynamics of these species.
How do beneficial alleles spread through these spatially-structured populations, and can we use spatial structure to detect them?
In Chapter III, I present my work investigating these questions, using a similar continuous-space simulation framework to uncover
a novel spatial signature of selective sweeps and then applying this observation to identify positively-selected alleles 
in a spatially-referenced genomic dataset of \textit{Anopheles gambiae} mosquitoes.


\chapter{Sweeps in Space: Leveraging Geographic Data to Identify Beneficial Alleles in \textit{Anopheles gambiae}}

\bigskip
This chapter was published in the journal \textit{Molecular Biology and Evolution} in 2025.
Scott T. Small, Peter L. Ralph, and Andrew D. Kern are co-authors on this work.
I conceptualized this study together with A.D.K. and P.L.R.; S.T.S. assisted in processing the 
\textit{Anopheles gambiae} genomic dataset and P.L.R. assisted with our conceptual model.
I performed all analyses with input from and wrote the manuscript with editorial assistance from all co-authors.

\bigskip
The citation for this publication is as follows:

\fullcite{rehmann2025sweeps}
\bigskip

\input{spatial_sweeps.tex}

\newpage
\section{Bridge}

As presented in Chapter III, spatial structure can be used to uncover adaptive variation at the population level, and indeed we successfuly identify
several candidate sweeping alleles in \textit{Anopheles gambiae}. 
However, \textit{Anopheles gambiae} is not the only malaria vector, and many \textit{Anopheles} species are segregating large amounts of potentially-adaptive variation.
In Chapter IV, I use a comparative approach to investigate the evolutionary forces shapping patterns of adaptive variation and divergence in 
three major \textit{Anopheles} malaria vectors.
This work reveals common selection pressures acting on these deeply-diverged vectors, with immune and metabolic pathways being particularly affected across a broad range of evolutionary timescales,
and highlights how local adaptation and diversifying selection maintain genetic variation in the same pathways.

\chapter{Convergent selection on nonsynonymous variation across deeply diverged \textit{Anopheles} vectors}

\bigskip
This chapter is in preparation for submission to the journal \textit{PLOS Genetics} and is coauthored with 
Andrew D. Kern. I conceptualized this study together with A.D.K. and performed all analyses; I wrote the manuscript 
with input and editorial assistance from A.D.K.

\newcommand{\comment}[1]{{\color{blue} #1}}
\newcommand{\ag}{Ag1000G }
\newcommand{\Anopheles}{\textit{Anopheles}}

\section{Introduction}

The \textit{Anopheles} mosquitoes that transmit human malaria
are among the most medically important insects on earth,
yet they remain remarkably difficult to control.
The genus is deeply diverged:
new world \textit{Anopheles} split from old world lineages
approximately 100 million years ago (MYA),
and more recent radiations across Africa and Asia
have given rise to a number of major vector species
and species complexes.
Recent population genomic studies have begun to characterize
the genetic diversity and evolutionary histories of these species
\citep{anopheles_comparative, darlingi, funestus, gambiae, minimus},
revealing a consistent pattern:
\textit{Anopheles} vectors are highly adaptable
and readily evolve to counter vector control efforts
\citep{gambiae, funestus}.
Understanding how these mosquitoes adapt so rapidly requires understanding
the evolutionary forces that shape patterns of genetic variation
across the genus.

Comparative genomic studies across the genus have revealed a striking capacity for genomic change:
\textit{Anopheles} genomes exhibit faster rates of gene gain and loss,
elevated gene shuffling on the X chromosome,
and rapid diversification of genes involved in chemosensation,
immunity, and insecticide detoxification \citep{anopheles_comparative}.
These patterns point to a genus whose genomes are unusually dynamic,
but they describe evolution between species
and leave open the question of what is happening within them.

While phylogenetic study of \textit{Anopheles} evolution has revealed the genus's evolvability
and highlighted the rapid rate at which their genomes adapt,
population genetics studies suggest that this evolvability may be due in part
to exceptionally high levels of segregating polymorphism.
In the African major vector \textit{Anopheles gambiae},
individual mosquitoes have 1.7 to 2.7 million variant sites
and 21\% of variants identified are triallelic;
in another major African vector \textit{An. funestus},
individuals have 2 to 3 million variant sites
and 15\% of variants have more than two alleles \citep{gambiae, funestus}.
Similar levels of diversity and proportions of multiallelic sites
have been reported in \textit{Anopheles minimus}, a major vector across Southeast Asia,
and \textit{An. darlingi}, the primary malaria vector in South America \citep{darlingi, minimus}.
Despite high amounts of global polymorphism,
these population genetic surveys have identified many instances of selection
acting on local populations,
particularly regarding adaptation to insecticides
\citep{Xue_Schrider_Kern_2021, Clarkson_Miles_Harding_2021, Barnes_Weedall_Ndula_2017},
and indeed while population structure in these species allows for local adaptation,
strong enough selection pressures can lead to the rapid spread of adaptive variants
such as insecticide resistance at the \textit{Vgsc} locus \citep{gambiae}.
This raises the question of what forces maintain polymorphism in these populations
while also allowing for rapid adaptation to local environmental pressures.
In most systems, nonsynonymous polymorphism is assumed to be mildly deleterious,
maintained by drift and purged inefficiently in small populations.
\textit{Anopheles} mosquitoes, however, have massive effective population sizes
spread across heterogeneous environments
with differential exposure to vector control measures across space.
We would therefore expect selection to act efficiently in these species,
both fixing beneficial alleles and removing deleterious ones.
Despite this, \textit{Anopheles} populations continue to segregate large amounts of nonsynonymous variation,
suggesting that some of this polymorphism may not be deleterious
but instead actively maintained by balancing selection or local adaptation.

What is missing is a comparative, gene-level understanding
of nonsynonymous variation across \textit{Anopheles} species ---
which genes are segregating functional polymorphism,
how that polymorphism relates to directional selection along lineages,
and whether similar selective pressures - in particular, vector control efforts -
 are producing convergent patterns
across deeply diverged vectors.
Here, we address these questions using two African species
(\textit{An. gambiae} and \textit{An. funestus})
and one Asian species (\textit{An. minimus}) (Figure \ref{fig:popdist}A).
The lineage leading to \textit{An. gambiae} diverged 30--50 MYA
and has more recently radiated into the \textit{gambiae} species complex,
while \textit{An. funestus} and \textit{An. minimus} are more closely related,
diverging approximately 22 MYA (Figure \ref{fig:popdist}B) \citep{anopheles_comparative}.
\textit{An. gambiae} and \textit{An. funestus} occupy similar ranges in sub-Saharan Africa
and are both highly anthropophilic;
they differ in larval site preference
but are jointly responsible for the majority of the continent's malaria burden,
transmitting primarily \textit{Plasmodium falciparum}.
\textit{An. minimus}, by contrast,
is a major vector in the Greater Mekong Subregion of Southeast Asia
with variable behavior and habitat preference;
it transmits primarily \textit{Plasmodium vivax} \citep{anophelesbook}.
All three species occupy heterogeneous environments
with variable insecticide exposure across both space and time.
Investigating both population genetic dynamics and
patterns of divergence, we identify parallel and convergent 
selection pressures that act on these highly-diverged vectors and shape their 
evolutionary trajectories.
We find that genes involved in immune response, insecticide resistance,
redox metabolism, cuticle formation, and sensory perception
both harbor excess nonsynonymous polymorphism within species
and are evolving under directional selection along lineages.
In \textit{An. gambiae} and \textit{An. funestus},
we further identify convergent spatial patterns of nonsynonymous differentiation
at two cytochrome P450 loci involved in insecticide resistance,
with parallel molecular changes in both species.
Together, these results suggest that local adaptation maintains
a reservoir of functional nonsynonymous variation in \textit{Anopheles},
providing substrate for rapid evolutionary responses
when selection pressures intensify or shift.


\begin{figure}[htp]
\includegraphics[width=\linewidth]{{figures/anopheles_tree.pdf}}
\centering
\caption{(A) Geographic distribution and (B) phylogeny of \textit{Anopheles gambiae}, \textit{funestus}, and \textit{minimus}.
}
\label{fig:popdist}
\end{figure}

\section{Methods}

\subsection{Genomic data \& processing} \label{sec:geno}

\textit{An. gambiae, funestus,} and \textit{minimus} population genetic datasets were obtained from the MalariaGen Vector Observatory.
All population VCFs were merged and then filtered (in the case of \textit{An. gambiae} and \textit{funestus}) using the provided site filters. 
Genome annotations were obtained in GFF format from VectorBase and processed using \texttt{gffutils 0.13}; reference genome and proteome sequences were obtained in Fasta format from VectorBase.
For all genes, the longest transcript was used for coding sequence analysis.

Prior to comparing patterns of nonsynonymous polymorphism and fixation between species, we identified one-to-one orthologs between \textit{An. gambiae, funestus}, and \textit{minimus} using \texttt{SonicParanoid 2.1.9} and aligned the protein sequence of all samples for a given ortholog using \texttt{MUSCLE 5.3}. Aligned amino acid sequences were mapped back to aligned nucleotide coding sequences using \texttt{Bio.Align.CodonAligner} from \texttt{BioPython 1.85}.

\subsection{Population genetics methods}

\subsubsection{Within species}

For every protein-coding gene in each species, we calculated within-species
 population genetics statistics considering the following sequences: 
the entire length of the gene, the coding sequence, the amino acid sequence, 
nonsynonymous sites ($N$), and synonymous sites ($S$). 
Tajima's $D$, nucleotide diversity ($\pi$), %Watterson's $\theta$, 
and Weir and Cockerham's $F_{ST}$ were calculated using a custom Python library, 
\texttt{WGSstats}, based on functions available in \texttt{scikit-allel} and \texttt{BioPython} 
(available at \comment{github URL}). 
For \textit{An. gambiae} and \textit{An. funestus}, statistics except for $F_{ST}$ 
were calculated on spatial cohort groups predetermined by the MalariaGEN data release 
\citep{gambiae, funestus} and then averaged across cohorts in order to minimize the impact of 
rare variants; for \textit{An. minimus} all samples were used at once.
To calculate $F_{ST}$ within species, we used the \texttt{HBDSCAN} algorithm implemented through \texttt{scikit-learn} 
spatially cluster sampling locations (with a minimum cluster size of five samples) and then
calculated Weir and Cockerham's $F_{ST}$ between all clusters for a given gene.
To assess the relationship between nonsynonymous and genic statistics,
 we fit linear models using the \texttt{stats} R package with the formula 
 $X_{N} \sim 0 + X_{G}$ where $X_{N}$ is the statistic for nonsynonymous sites 
 sand $X_{G}$ is the statistic for the entire gene \citep{Rstats}.
We obtained $p$ values for each statistic based on their relationship to the 
overall within-species distribution and corrected for FDR using the Benjamini-Hochberg method. 
Outlier genes were then identified as genes with $p < 0.05$ for $D_N$, $\pi_N$, $F_{ST_{N}}$, 
and as genes with ${p_{N}}/{p_{S}} > 1$ for nonsynonymous to synonymous polymorphism ratios.
To test for GO term enrichment amongst outlier genes, performed gene set enrichment analysis 
using R package \texttt{topGO} \citep{RtopGO} and GO term 
annotations obtained from VectorBase \citep{vectorbase}.
To test for the statistical significance of overlap between lists of outlier genes, 
we used R package \texttt{SuperExactTest} \citep{Rset}. 

\subsubsection{Between species} \label{sec:between}

We performed polarized gene-wise McDonald-Kreitman tests on each one-to-one 
ortholog identified in \nameref{sec:geno} using Python package \texttt{mkado} \citep{mkado}. 
In order to reduce the influence of rare variants on our analysis, we used a subset 
of samples from each species selected from the same geographic region and time to conduct the test. 
For \textit{An. gambiae} we used the AG1000G-GN-A cohort, for \textit{funestus} we 
used set 1230-VO-GA-CF-AYALA-VMF00045, and for \textit{minimus} we used samples 
taken from coordinate location $(13.667, 104.982)$ in 2014.
For \textit{An. gambiae}, we used \textit{funestus} as an outgroup and polarized 
using \textit{An. minimus}; for \textit{An. funestus} we used \textit{gambiae} as an 
outgroup and polarized using \textit{An. minimus}; for \textit{An. minimus} we used
 \textit{funestus} as an outgroup and polarized using \textit{An. gambiae}.
With \texttt{mkado}, we also performed asymptotic genome-wide McDonald-Kreitman 
tests using the same ingroup-outgroup-polarization pairings.

To test for positive selection along the \textit{gambiae}, \textit{funestus}, and \textit{minimus} 
branches, we performed multi-ratio branch model tests on each one-to-one ortholog with the 
focal species as the foreground branch using \texttt{crabML} \citep{crabml2025, yang2007paml}. 
In order to account for the high degree of polymorphism within species, we ran ten test 
replicates for each gene, using randomly-chosen samples from each species as representative inputs. 

\subsubsection{$F_{ST}$ correlation} \label{sec:fst}

To investigate spatial correlation between orthologous nonsynonymous $F_{ST}$ outlier genes in \textit{An. gambiae} and \textit{funestus}, 
we calculated pairwise Hudson's $F_{ST_{N}}$ (within each species) between clusters of sampling points for each gene.
Sample clusters were identified using an HDBSCAN algorithm implemented through \texttt{scikit-learn} with a minimum cluster size of five individuals.
To identify the average spatial axis of genetic differentiation for a given gene, we then treated each cluster pair as a west$\rightarrow$east vector and averaged these vectors, weighted by their $F_{ST_{N}}$:
$$\bar{(x_{F_{ST_{N}}}, y_{F_{ST_NS}}}) =\frac{\sum_{i=1}^{n}{F_{ST_{NS_i}} (x_{i1}, y_{i1}, x_{i2}, y_{i2})}}{\sum_{i=1}^{n}{F_{ST_{NS_i}}}}$$
We then tested for spatial similarity between \textit{An. gambiae} and \textit{funestus} differentiation using the cosine similarity between average $F_{ST}$ vectors for a given pair of orthologs.
To test for significance, we randomly permuted the $F_{ST_{N}}$ values used to calculate the average $F_{ST_{N}}$ vector ($N=1000$) and compared observed and permuted cosine similarities.

For genes with significantly similar nonsynonymous $F_{ST}$,
 we probed which mutations were driving west$\rightarrow$east differentiation by 
 calculating $F_{ST}$ between western and eastern sampling locations for each nonsynonymous mutation. 
 Mutations with $F_{ST}$ greater than 0.2 were considered driving 
 mutations and retained for further analysis. To identify convergent driving 
 mutations between \textit{An. gambiae} and \textit{funestus}, 
 we analyzed the spatial distribution of each variant and used 
 \texttt{ThermoMPNN} to assess the impacts of mutations on protein stability \citep{thermompnn}.

\section{Results}

\subsection{Nonsynonymous polymorphism is abundant and enriched for immune, sensory, and metabolic functions}

To characterize nonsynonymous variation within each species,
%we first examined genome-wide patterns of polymorphism, site frequency spectra,
we examined patterns of polymorphism, site frequency spectra, 
and spatial differentiation at nonsynonymous sites
relative to genic and synonymous baselines.
Reported genome-wide nucleotide diversity is high in all three species:
$\pi =1.5\times10^{-2}$  in \textit{An. gambiae}, between
$1.4\times10^{-2}$ and $1.5\times10^{-2}$ in \textit{An. funestus},
and between $1.3\times10^{-2}$ and $1.4\times10^{-2}$ in \textit{An. minimus},
consistent with their large effective population sizes \citep{gambiae, funestus, minimus}.
In our analysis, genic $\pi$ ($\pi_G$) is approximately equal to genome-wide $\pi$ in each species --- 
except for in \textit{An. gambiae}, where it is slightly higher ---
and nonsynonymous $\pi$ ($\pi_{N}$) is one order of magnitude lower than $\pi_G$
($\pi_{N} =2.58\times10^{-3}$, $4.17\times10^{-3}$,
and $4.10\times10^{-3}$ for \textit{gambiae}, \textit{funestus},
and \textit{minimus}, respectively) (Figure \ref{fig:popstats}, Table \ref{table:statstable}).
For comparison, genome-wide $\pi$ in African \textit{D. melanogaster}
is $7.5\times10^{-3}$,
with $\pi_N = 1.5\times10^{-3}$ and $\pi_S = 1.9\times10^{-2}$ \citep{drosophila}.
Nonsynonymous $\pi$ in \textit{An. gambiae} is thus almost double that of
African \textit{D. melanogaster};
despite \textit{An. gambiae}'s substantially higher genic diversity,
in \textit{An. funestus} and \textit{minimus}
$\pi_{N}$ is higher still.
Nonsynonymous $\pi$ is significantly correlated with gene-wide and synonymous $\pi$
in all three species (Figure \ref{fig:pilm}).
However, the relationship between nonsynonymous and genic polymorphism differs among species:
\textit{An. gambiae} appears more constrained,
while \textit{An. funestus} and \textit{An. minimus} show lower genic $\pi$
but higher nonsynonymous $\pi$.

Turning to the site frequency spectrum,
almost all genes in each species have negative Tajima's $D$ values
(median $D_G =$ -1.48, -1.38, -1.83
and median $D_{N} =$ -1.97, -1.88, -2.27
for \textit{gambiae}, \textit{funestus}, and \textit{minimus}, respectively),
consistent with recent population expansions.
Negative genome-wide Tajima's $D$ is also observed in \textit{D. melanogaster},
where it is driven primarily by linked selection
rather than demography \citep{drosophila}.
Although across the genome the distribution of nonsynonymous Tajima's $D$
is shifted below genic Tajima's $D$ in each species,
linear models of the per-gene relationship between nonsynonymous and genic $D$
reveals that individual genes tend to have higher nonsynonymous than genic $D$
(coefficient $\beta = 1.227, 1.237, 1.173$ for \textit{gambiae, funestus}, and \textit{minimus};
$p < 2\times10^{-16}$ and $\rho > 0.94$ for all three models;
Figure \ref{fig:dlm}).
This pattern --- nonsynonymous $D$ shifted upward relative to genic $D$
on a per-gene basis --- is consistent with balancing forces
maintaining nonsynonymous variation at a subset of genes.

Finally, genome-wide distributions of genic and nonsynonymous $F_{ST}$ are nearly identical
within each species, with a close to 1:1 relationship between the two measures
(Figures \ref{fig:popstats}, \ref{fig:fstlm}).
In most organisms, purifying selection would be expected to depress nonsynonymous $F_{ST}$
relative to genic $F_{ST}$.
The absence of this pattern here suggests that local adaptation
may be elevating nonsynonymous differentiation,
consistent with the high levels of nonsynonymous polymorphism
we observe across species.

\begin{table}[htp]
\centering
\begin{tabular}{@{}lllllll@{}}
\toprule
\textbf{Species}      & \textbf{$\pi_{G}$} & \textbf{$\pi_N$} & \textbf{Tajima's $D_G$} & \textbf{Tajima's $D_N$} & \textbf{$F_{ST_G}$} & \textbf{$F_{ST_{N}}$} \\ \midrule
\textit{An. gambiae}  & 0.02348            & 0.00258             & -1.48478                & -1.96957                   & 0.04450             & 0.03859                \\
\textit{An. funestus} & 0.01467            & 0.00417             & -1.38169                & -1.88379                   & 0.08568             & 0.07314                \\
\textit{An. minimus}  & 0.01206            & 0.00408             & -1.82551                & -2.26612                   & 0.01562             & 0.01224                \\ \bottomrule
\end{tabular}
\caption{Summary of genic and nonsynonymous population genetic statistics for 
\textit{An. gambiae}, \textit{funestus}, and \textit{minimus}. Tajima's $D$ values are reported as the median across genes; other 
values are reported as the mean across genes.}
\label{table:statstable}
\end{table}


Taken together, the $\pi$, Tajima's $D$, and $F_{ST}$ results point to a shared pattern
across all three species: nonsynonymous variation is abundant,
shows signatures of balancing selection on a per-gene basis,
and is not depleted by purifying selection to the extent seen in other organisms.
We next examine which genes and functional categories
drive these genome-wide patterns in each species.

\begin{figure}[htp]
\includegraphics[width=\linewidth]{{figures/population_metrics_violin_plots.pdf}}
\centering
\caption{Distribution of genic (lighter color) and nonsynonymous (darker color) $\pi$, Tajima's $D$, and $F_{ST}$ values within \textit{An. gambiae}, \textit{funestus}, and \textit{minimus}.
}
\label{fig:popstats}
\end{figure}

In \textit{An. gambiae}, genes with elevated nonsynonymous polymorphism
include cytochrome P450s, which are involved in insecticide resistance,
immune genes such as CLIP-domain serine protease (AGAP011790),
and cuticular genes such as chitinase (AGAP006191).
Among high-$\pi_{N}$ genes, we find significant enrichment
for sensory perception of smell, stimulus response,
stress response [DNA damage response, cellular response to stress],
proteolysis, and proteoglycan metabolism
(Figure \ref{fig:goenrich}, Table \ref{gambGOpi}).
The overrepresentation of these functional groups
likely reflects varying environmental pressures
across \textit{An. gambiae}'s range,
which covers the majority of sub-Saharan Africa:
sensory perception and stimulus response may involve
host-seeking behavior or insecticide avoidance,
both of which vary across the landscape,
while stress response is also associated with insecticide resistance.
The enrichment of proteoglycan metabolism genes may reflect balancing selection
on an immune trait, as proteoglycans mediate \textit{Plasmodium} interactions
\citep{Armistead_Wilson_vanKuppevelt_Dinglasan_2011}.

Genes with elevated nonsynonymous Tajima's $D$ echo several of these categories,
with enrichment for sensory perception of smell and proteolysis,
a key component of blood meal digestion,
as well as glutathione, sulfur compound, and NADP metabolic processes
(Table \ref{gambGOd}).
These metabolic categories are all involved in maintaining
redox equilibrium \citep{redox_fecundity}.
The production of reactive oxygen species
is a key part of \textit{Anopheles} immune defense,
and redox states have been shown to impact fitness in \textit{An. gambiae}
\citep{redox_fecundity, redox_gutmicro, redox_plasmodium}.
The convergence of $\pi_{N}$ and Tajima's $D_{N}$ enrichments
on overlapping functional categories suggests that balancing selection
and local adaptation help maintain functional polymorphism
in \textit{An. gambiae} redox metabolism,
with potential downstream impacts on immunity, fecundity, and insecticide susceptibility.

A similar pattern emerges in \textit{An. funestus},
where genes with elevated nonsynonymous polymorphism
include a number of immune genes ---
we observe 13 CD109 antigen-like genes, leucine-rich repeat proteins,
and a TsetseEP-like gene among outliers.
Among high-$\pi_{N}$ genes, we find significant enrichment
for reproduction, development, and protein metabolism
(Figure \ref{fig:goenrich}, Table \ref{afunGOpi}).

Genes with elevated nonsynonymous Tajima's $D$ also include immune genes,
as well as multiple phenoloxidase-activating enzymes.
GO enrichment for high-$D_{N}$ genes reveals protein metabolic processes
along with a number of catabolic and transport processes
(Table \ref{afunGOd}).
Notably, metabolic enrichments center on nucleoside diphosphate
and carbohydrate metabolism,
particularly purine ribonucleotide diphosphates,
which are involved in the pentose phosphate pathway (PPP),
a critical component of \textit{Anopheles} redox metabolism
\citep{Kanehisa_Furumichi_Sato_Matsuura_Ishiguro-Watanabe_2024}.
The presence of phenoloxidase-activating enzymes among high-$D_{N}$ genes
reinforces this redox connection,
as phenoloxidase activity is tightly linked to reactive oxygen species production.
As in \textit{An. gambiae}, the overlap between $\pi_{N}$ and Tajima's $D_{N}$
enrichments on metabolic processes suggests that balancing selection
and spatially varying pressures maintain functional polymorphism
in \textit{An. funestus},
with protein and carbohydrate metabolism genes
likely subject to varying dietary and immunological demands across the species range.

The picture differs in \textit{An. minimus}.
Genes with elevated nonsynonymous polymorphism
include cuticular proteins and immune genes such as a CLIP-domain serine protease,
and elevated nonsynonymous Tajima's $D$ flags additional immune genes.
Among high-$\pi_{N}$ genes, we find enrichment primarily
for broad regulatory functions
(Figure \ref{fig:goenrich}, Table \ref{aminGOpi}).
High-$D_N$ genes show a similar regulatory signature,
with enrichment for regulation of catalytic activity,
protein metabolic process, and macromolecule metabolic process,
as well as sensory perception of chemical stimulus
(Table \ref{aminGOd}).
The sensory enrichment parallels the pattern observed in \textit{An. gambiae},
but the overall enrichment profile in \textit{An. minimus}
is notably distinct from the other two species:
where \textit{An. gambiae} and \textit{An. funestus} show enrichment
for specific metabolic and immune pathways,
\textit{An. minimus} enrichments center on higher-level regulatory processes,
suggesting that nonsynonymous variation in this species
may be maintained through different selective pressures
or different evolutionary responses to similar pressures.

Across all three species, immune genes and sensory genes
consistently appear among nonsynonymous outliers,
while the metabolic signatures of balancing selection ---
redox metabolism in \textit{An. gambiae} and \textit{An. funestus},
broad regulatory processes in \textit{An. minimus} ---
suggest both shared and lineage-specific pressures
shaping nonsynonymous variation in these vectors.

\begin{figure}[htp]
\includegraphics[width=\linewidth]{{figures/population_metrics_go_enrichment_ns_pi.pdf}}
\centering
\caption{Significant GO enrichments for nonsynonymous $\pi$ outliers in \textit{An. gambiae}, \textit{funestus}, and \textit{minimus}.
}
\label{fig:goenrich}
\end{figure}

\subsection{Adaptive divergence targets overlapping functional categories across lineages}

\subsubsection{McDonald Kreitman test}

Having established that nonsynonymous variation is abundant
and bears signatures of balancing selection,
we next asked whether comparisons of polymorphism and divergence
show evidence of non-neutral evolution. 
Using the asymptotic McDonald-Kreitman test \citep{asym_mk},
we estimated the proportion of amino acid fixations driven by directional selection ($\alpha$)
for all one-to-one orthologs across the three species.
Despite the high levels of nonsynonymous polymorphism documented above,
genome-wide $\alpha$ exceeds 0.60 in all three lineages,
with the lowest value in \textit{An. gambiae} ($\alpha = 0.600$)
and the highest in \textit{An. funestus} ($\alpha = 0.655$).
These values are comparable to or slightly higher than
the asymptotic $\alpha$ of 0.571 estimated for \textit{D. melanogaster}
\citep{asym_mk, Haller_Messer_2017},
indicating similarly pervasive adaptive protein evolution
despite very different ecology and life history.
However, the number of individual genes
with significant evidence of directional selection is small,
indicating that while strong positive selection acts on a subset of genes,
the majority continue to segregate nonsynonymous variants
(Figure \ref{fig:mktest}).

\begin{figure}[htp]
\includegraphics[width=\linewidth]{{figures/polarized_mk_tests_all_genes.pdf}}
\centering
\caption{Results of the asymptotic McDonald-Kreitman test for all one-to-one orthologs across \textit{An. gambiae}, \textit{An. funestus}, and \textit{An. minimus}.
Light blue points represent genes with significant evidence of adaptive divergence after Boneferroni correction ($p < 0.05$), 
while dark blue points represent genes with a significant excess of nonsynonymous polymorphism.}
\label{fig:mktest}
\end{figure}

In \textit{An. gambiae}, 1238 of 5149 one-to-one orthologs (24\%)
show evidence of non-neutral evolution ($p < 0.05$),
but only 171 of these have $\alpha$ values
indicating adaptive directional selection ($\alpha > 0.5$).
The gene with the highest significant $\alpha$,
fatty acyl-CoA reductase 2 (AGAP004784),
is involved in cuticle formation and has an estimated 100\%
of its nonsynonymous fixations due to positive selection
($p = 3.406 \times 10^{-3}$).
Several sensory genes also show strong directional selection,
including odorant-binding protein 11
(AGAP002025, $\alpha = 0.900$, $p = 6.926\times10^{-3}$),
sensory appendage protein 1
(AGAP008051, $\alpha = 0.816$, $p = 2.887\times10^{-2}$),
and odorant-binding protein 69
(AGAP006368, $\alpha = 0.814$, $p = 1.393\times10^{-2}$).
Notably, sensory perception of smell is also enriched
among high-$\pi_N$ genes (see above),
indicating that this functional category potentially experiences both balancing
and directional selection on different loci.
Immune genes with elevated $\alpha$ include a mannose-binding C-type lectin
(AGAP029559, $\alpha = 0.818$, $p = 1.449\times10^{-12}$),
a CLIP-domain serine protease
(AGAP011789, $\alpha = 0.681$, $p = 2.162\times10^{-2}$),
and a leucine-rich immune protein
(AGAP006348, $\alpha = 0.616$, $p = 1.044\times10^{-3}$).
One of the strongest signals of adaptive divergence
is in a heparan-sulfate 6-O-sulfotransferase
(AGAP001444, $\alpha = 0.937$, $p = 1.195\times10^{-2}$).
Heparan sulfate mediates \textit{Anopheles}--\textit{Plasmodium} interactions,
and this signal may reflect coevolution between vector and parasite
\citep{Armistead_Wilson_vanKuppevelt_Dinglasan_2011,
Sinnis_Coppi_Toida_Toyoda_Kinoshita-Toyoda_Xie_Kemp_Linhardt_2007}.

\begin{table}[htp]
\centering
\begin{tabular}{@{}llll@{}}
\toprule
\textbf{Gene} & \textbf{$p$ value}    & \textbf{$\alpha$} & \textbf{Product Description}                                                                      \\ \midrule
AGAP004784    & $3.406\times10^{-3}$  & 1.00                 & fatty acyl-CoA reductase 2                                                                        \\
AGAP009154    & $7.697\times10^{-5}$  & 0.96                 & actin-related protein 2/3 complex subunit 2                                                       \\
AGAP001444    & $1.195\times10^{-2}$  & 0.94                 & Heparan-sulfate 6-O-sulfotransferase                                                              \\
AGAP005400    & $4.334\times10^{-7}$  & 0.92                 & four and a half LIM domains                                                                       \\
AGAP001903    & $1.662\times10^{-17}$ & 0.91                 & NAD-dependent malate dehydrogenase                                                                \\
AGAP005696    & $1.664\times10^{-3}$  & 0.91                 & cuticular protein from two-cysteine family 1                                                      \\
AGAP005991    & $5.071\times10^{-3}$  & 0.91                 & 60S ribosomal protein L14                                                                         \\
AGAP007668    & $3.841\times10^{-2}$  & 0.91                 & Eukaryotic translation initiation factor 3 subunit G                                              \\
AGAP002025    & $6.926\times10^{-3}$  & 0.90                 & odorant-binding protein 11                                                                        \\
AGAP005958    & $4.641\times10^{-4}$  & 0.86                 & Yellow-g                                                                                          \\ \bottomrule
\end{tabular}
\caption{Top 10 genes with significant evidence of adaptive divergence in \textit{An. gambiae}}
\label{table:gambMK}
\end{table}

Testing for functional enrichment among high-$\alpha$ genes,
we find enrichment for translation, protein-containing complex assembly,
and nuclear transport --- categories likely reflecting general species divergence,
though changes in protein complex assembly could also involve
complement-like complexes central to \textit{Anopheles} immune response
\citep{Povelones_Upton_Sala_Christophides_2011,
Povelones_Waterhouse_Kafatos_Christophides_2009}.
We also find enrichment for stress response,
carbohydrate derivative catabolism, and organophosphate catabolism.
As with sensory genes, stress response genes
appear to experience both balancing selection (excess nonsynonymous polymorphism)
and directional selection (elevated $\alpha$) on different loci.
The enrichment for organophosphate catabolism is particularly striking
and may represent a direct adaptation
to the phosphate-containing insecticides deployed in vector control.

A similar proportion of genes show non-neutral evolution in \textit{An. funestus}
(1249 of 5149, 24.2\%),
though a notably larger fraction --- 710 genes --- have $\alpha$ values
indicating adaptive divergence, compared to 171 in \textit{An. gambiae}.
Cuticle formation, sensory perception, insecticide resistance,
and immune response are again common themes among high-$\alpha$ genes.
Fourteen genes in \textit{An. funestus} lack any nonsynonymous polymorphism
despite nonsynonymous divergence ($\alpha = 1.0$),
including Tango10 (AGAP000791, $p = 2.909\times10^{-5}$),
a circadian rhythm protein,
and sensory appendage protein 1
(AGAP008051, $p = 3.472 \times 10^{-4}$),
which also shows elevated $\alpha$ in the \textit{gambiae} lineage,
indicating that this gene is diverging under directional selection
in both species independently.
Yellow-g (AGAP005958) provides another case of parallel divergence,
with $\alpha = 0.873$ ($p = 1.804\times10^{-4}$) in \textit{An. gambiae}
and $\alpha = 1.0$ ($p = 7.884\times10^{-5}$) in \textit{An. funestus}.
Its ortholog inhibits \textit{Plasmodium vivax} in \textit{An. dirus}
\citep{yellowg_plasmodium}
and prevents egg desiccation in \textit{Aedes} mosquitoes
\citep{yellowg_desiccation},
either of which could explain its adaptive divergence along both lineages.

\begin{table}[htp]
\centering
\begin{tabular}{@{}llll@{}}
\toprule
\textbf{Gene} & \textbf{$p$ value}   & \textbf{$\alpha$} & \textbf{Product Description}                         \\ \midrule
AGAP000308    & $4.426\times10^{-2}$   & 1.00                 & proteasome activator subunit 3 (PA28 gamma)          \\
AGAP000791    & $2.909\times10^{-5}$ & 1.00                 & Tango10                                              \\
AGAP001515    & $3.386\times10^{-2}$   & 1.00                 & unspecified product                                  \\
AGAP001805    & $9.179\times10^{-3}$  & 1.00                 & 60S ribosomal protein L13                            \\
AGAP002053    & $1.224\times10^{-2}$   & 1.00                 & unspecified product                                  \\
AGAP002935    & $1.312\times10^{-2}$   & 1.00                 & Eukaryotic translation initiation factor 3 subunit F \\
AGAP002936    & $4.021\times10^{-2}$   & 1.00                 & unspecified product                                  \\
AGAP003462    & $2.819\times10^{-2}$   & 1.00                 & protein arginine N-methyltransferase 1               \\
AGAP005991    & $4.220\times10^{-3}$  & 1.00                 & 60S ribosomal protein L14                            \\
AGAP005958    & $7.884\times10^{-5}$ & 1.00                 & Yellow-g                                             \\ \bottomrule
\end{tabular}
\caption{Top 10 genes with significant evidence of adaptive divergence in \textit{An. funestus}}
\label{table:afunMK}
\end{table}

Functional enrichment among high-$\alpha$ genes in \textit{An. funestus}
reveals a broader set of categories than in \textit{An. gambiae},
including reproduction, cell cycle, nuclear transport,
translational initiation, development, sensory perception, and locomotion.
Reproductive and developmental enrichments
likely reflect the rapid evolution characteristic of these functional categories,
and nuclear transport and translation may be attributed to general species divergence.
More striking are the enrichments for locomotion and signal transduction,
both of which may be associated with adaptation to insecticide exposure:
pyrethroids in particular cause locomotor deficits,
and resistance is associated with changes in locomotion
\citep{pyretheroid-loc, pyretheroid-fit},
while signal transduction is disrupted by several classes of insecticides
\citep{signal_transduction_insecticides}
and may also contribute to behavioral resistance via avoidant behavior.


The \textit{An. minimus} lineage shows markedly less non-neutral evolution:
only 815 of 5149 orthologs (15.8\%) reject neutrality,
and just 215 have $\alpha$ values indicating adaptive divergence.
As in \textit{An. funestus}, a number of genes lack any nonsynonymous polymorphism
despite nonsynonymous divergence ($\alpha = 1.0$),
including two post-translational modification enzymes:
serine/threonine protein phosphatase 6
(AMIN003725, $p = 1.493\times10^{-2}$)
and protein arginine N-methyltransferase 1
(AMIN000933, $p = 5.432\times10^{-3}$).
Other high-$\alpha$ genes include immune proteins
such as spaetzle-like cytokine 1
(AMIN005499, $\alpha = 0.808$, $p = 9.834\times10^{-5}$)
and a mannose-binding C-type lectin
(AMIN016218, $\alpha = 0.756$, $p = 3.087\times10^{-2}$),
as well as cuticular proteins such as fatty acyl-CoA reductase (AGAP001084)
and sensory proteins such as odorant-binding protein 40 (AGAP002191).
Although we do not find evidence for adaptive divergence in Yellow-g
as we do in \textit{An. gambiae} and \textit{An. funestus},
a paralog, Yellow-g2, is adaptively diverging in \textit{An. minimus}
(AMIN010972, $\alpha = 0.796$, $p = 4.141\times10^{-3}$).
Yellow-g2's ortholog in \textit{Aedes} prevents egg desiccation
\citep{yellowg_desiccation},
and it may play a role analogous to Yellow-g
in \textit{Plasmodium} inhibition \citep{yellowg_plasmodium},
suggesting that selection on this gene family
is a recurring theme across \textit{Anopheles} lineages.

\begin{table}[htp]
\centering
\begin{tabular}{@{}llll@{}}
\toprule
\textbf{Gene} & \textbf{$p$ value}     & \textbf{$\alpha$} & \textbf{Product Description}                                       \\ \midrule
AGAP002762    & $1.975\times10^{-2}$   & 1.00                 & Chaperonin containing TCP1 subunit 7                               \\
AGAP003017    & $1.493\times10^{-2}$   & 1.00                 & serine/threonine-protein phosphatase 6 catalytic subunit           \\
AGAP003462    & $5.432\times10^{-3}$   & 1.00                 & protein arginine N-methyltransferase 1                             \\
AGAP008041    & $3.147\times10^{-3}$   & 1.00                 & RING finger protein 121                                            \\
AGAP009863    & $2.319\times10^{-3}$   & 1.00                 & eukaryotic translation initiation factor 4A, isoform 2             \\
AGAP011706    & $2.348\times10^{-2}$   & 1.00                 & 60S ribosomal protein L27a                                         \\
AGAP011937    & $8.950\times10^{-4}$   & 1.00                 & unspecified product                                                \\
AGAP001805    & $1.942\times10^{-3}$   & 0.96                 & 60S ribosomal protein L13  \\
AGAP004819    & $4.690\times10^{-3}$   & 0.95                 & receptor expression-enhancing protein 5                            \\
AGAP005542    & $2.030\times10^{-9}$   & 0.94                 & Queuine tRNA-ribosyltransferase subunit QTRTD1 homolog             \\ \bottomrule
\end{tabular}
\caption{Top 10 genes with significant evidence of adaptive divergence in \textit{An. minimus}}
\label{table:aminMK}
\end{table}

Functional enrichment among high-$\alpha$ genes in \textit{An. minimus}
highlights programmed cell death, translation, nuclear transport,
purine nucleotide metabolism, and phosphorus metabolism.
As in the other two species, translation and nuclear transport
likely reflect general species divergence.
Programmed cell death is a hallmark of both development and immune response;
its enrichment among adaptively diverging genes
may reflect directional selection on immune function.
The enrichment of purine ribonucleotide and ribose phosphate metabolism
parallels the polymorphism results in \textit{An. funestus}
and may signal functional divergence in redox metabolism,
driven by selection on immune defense or insecticide detoxification.
Finally, the enrichment for phosphorus metabolism
echoes the organophosphate catabolism signal in \textit{An. gambiae}
and likely reflects adaptation to organophosphate insecticides
deployed across these species' ranges.

Across all three species, the McDonald-Kreitman results reveal a consistent picture:
genome-wide $\alpha$ values comparable to \textit{Drosophila}
indicate pervasive adaptive protein evolution,
yet the genes individually under directional selection
are enriched for the same functional themes ---
immune response, sensory perception, cuticle formation,
and insecticide-related metabolism ---
that also harbor excess nonsynonymous polymorphism.
Several genes, including sensory appendage protein 1 and the Yellow-g family,
show parallel adaptive divergence across lineages,
reinforcing the convergence of selective pressures
acting on these deeply diverged vectors.

%\subsection{Balancing selection maintains nonsynonymous polymorphism}

%While we observe a suite of genes evolving under directional selection using the McDonald Kreitman test, \comment{there are also a lot of significant hits on genes with an excess of amino acid polymoprphism (\textit{ie}, a neutrality index greater than one}. Further, while we observe a strong negative relationship between the neutrality index and overall genic Tajima's $D$ in significant genes across all three species, there is a positive relationship between the neutrality index and nonsynonymous Tajima's $D$. Though the vast majority of genes in our analysis have negative Tajima's $D$ values (both genic and nonsynonymous), the correlation between excess amino acid polymorphism and increased $D_N$ suggests that this polymorphism is maintained in part due to balancing selection.
%Across all three species tested, genes with a significant excess of amino acid polymorphism are enriched for tRNA and RNA processing; functional variation in these processes may give rise to and maintain phenotypic variation that is not identifiable at the genotypic level.



\subsubsection{Phylogenetic testing}

To complement the McDonald-Kreitman analysis,
we used a branch model \citep{yang2007paml} to test for positive selection
($\omega > 1$) along the \textit{gambiae}, \textit{funestus},
and \textit{minimus} branches.
Because population-level samples contain segregating variation
that can influence branch-model inference,
we ran the test ten times per ortholog
using alignments of randomly drawn haplotypes from each species.
This design distinguishes genes where all sampled haplotypes
show $\omega > 1$ --- indicating fixation of adaptive substitutions ---
from genes where only some haplotypes do,
suggesting that adaptive divergence is still segregating within the species (\nameref{sec:between}).

In \textit{An. gambiae}, genes with significant $\omega > 1$
are enriched for signal transduction,
and we observe elevated $\omega$ on segregating haplotypes
in immune genes such as LRRCT domain-containing protein
and in cuticular proteins such as fatty acyl-CoA reductase.
In \textit{An. funestus}, enrichments shift to translation,
mitosis, and polyamine metabolism.
LRRCT domain-containing protein again shows elevated $\omega$
across all sampled haplotypes --- a parallel with \textit{An. gambiae} ---
while other immune genes such as mannose-binding C-type lectin
show $\omega > 1$ on only a subset of haplotypes,
illustrating how adaptive divergence can coexist
with segregating nonsynonymous variation at the same locus.
In \textit{An. minimus}, enrichments center on ATP metabolism,
chitin metabolism, amino sugar metabolism,
and purine ribonucleotide metabolism ---
the last of which echoes the purine metabolic enrichments
observed in both the polymorphism and MK analyses.
Unlike \textit{An. gambiae} and \textit{An. funestus},
immune and cuticular genes in \textit{An. minimus}
do not show haplotype-level variation in $\omega$,
consistent with the more generic regulatory enrichments
observed among \textit{minimus} nonsynonymous polymorphism outliers.

\subsection{Convergent nonsynonymous spatial differentiation between \textit{An. gambiae} and \textit{funestus}}

The analyses above show that \textit{An. gambiae} and \textit{An. funestus}
share many of the same functional targets of selection.
Because these two species are broadly sympatric across sub-Saharan Africa,
we next asked whether nonsynonymous spatial differentiation
within each species follows similar geographic patterns.
Among the one-to-one orthologs,
seven exhibit elevated nonsynonymous $F_{ST}$ in both species,
including three cytochrome P450s, a cholesterol transporter,
and a Niemann-Pick Type C-2 protein.
To test whether differentiation in these genes follows similar spatial axes,
we averaged nonsynonymous $F_{ST}$ vectors across sampling locations
for each gene and compared the resulting vectors between species
(\nameref{sec:fst}).
Three genes share significantly similar axes of nonsynonymous differentiation:
two cytochrome P450s at the Cyp6p locus ---
Cyp6p5 (AGAP002866) and Cyp6p2 (AGAP002869) ---
and protein FAM91A1 (AGAP002874) (Figure \ref{fig:fst_corr}).
The Cyp6p locus is a cluster of cytochrome P450 genes
whose mutations are implicated in insecticide resistance,
and the shared east-west axis of differentiation across species
suggests convergent adaptation to spatially varying insecticide exposure.
FAM91A1 has no well-documented role in \textit{Anopheles} biology,
but its convergent spatial differentiation with the Cyp6p genes
may indicate a related function in insecticide response.

In Cyp6p5, we identify four nonsynonymous mutations
contributing to west-east $F_{ST_{N}}$ in \textit{An. gambiae} (AGAP002866)
and seven in \textit{An. funestus} (AFUN2\_013651).
Of these, three in each species share the same spatial constraint,
occurring almost exclusively in the northeast of the sampling region
in the Central African Republic, Uganda, and northwest 
Tanzania (Figure \ref{fig:cyp6p5}A).
In both species, the spatially constrained mutations cluster
near the N-terminus or channel entrance of the protein,
suggesting convergent molecular evolution in this region
(Figure \ref{fig:cyp6p5}B).
We used ThermoMPNN to estimate the impact of individual mutations
on protein folding stability
(though current methods cannot model the effects of co-occurring mutations).
In the \textit{gambiae} Cyp6p5, N343T and V87I are moderately destabilizing
($\Delta\Delta G = 0.243$ and $0.083$ kcal/mol, respectively),
and H36C, while only slightly destabilizing ($\Delta\Delta G = 0.038$ kcal/mol),
substitutes a sulfur-containing cysteine for a histidine
near the protein's first $\alpha$ helix,
with potentially outsized impacts on folding or catalytic activity.
All three mutations segregate on the same haplotype.
In the \textit{funestus} Cyp6p5, V14L and R169H are moderately stabilizing
($\Delta\Delta G = -0.136$ and $-0.022$ kcal/mol),
while P171Q is moderately destabilizing ($\Delta\Delta G = 0.248$ kcal/mol),
likely due to substitution of a proline in a hairpin turn
between two $\alpha$ helices.
R169H and P171Q segregate on the same haplotype,
while V14L does not co-occur with either.
The remaining four \textit{funestus} mutations
are spatially constrained to the far southeast or northwest
and may represent independent adaptive responses
to insecticide exposure (Figure \ref{fig:afun_cyp}).
Although individual mutation effects on protein stability are moderate,
co-occurrence on the same haplotype in both species
raises the possibility of non-additive epistatic effects on protein function.


\begin{figure}[htp]
\includegraphics[width=0.9\linewidth]{{figures/AGAP002866_mutations.pdf}}
\centering
\caption{(A) Frequency of spatially-constrained nonsynonymous Cyp6p5 mutations in \textit{An. gambiae} (left) and \textit{funestus} (right). (B) Position of spatially-constrained nonsynonymous mutations on the Cyp6p5 protein in \textit{An. gambiae} (left) and \textit{funestus} (right). Arrow indicates channel entrance.
}
\label{fig:cyp6p5}
\end{figure}

A similar pattern emerges in Cyp6p2.
We identify four nonsynonymous mutations
contributing to west-east $F_{ST_{N}}$ in \textit{An. gambiae} (AGAP002869)
and nine in its \textit{An. funestus} ortholog.
One mutation in each species --- P449S in \textit{gambiae}
and Y391D in \textit{funestus} ---
shares the same spatial constraint
and maps to the same region of the protein structure
(Figure \ref{fig:cyp6p2}).
Both are strongly destabilizing:
P449S ($\Delta\Delta G = 1.067$ kcal/mol) disrupts a hairpin turn
near the heme-binding pocket,
and Y391D ($\Delta\Delta G = 2.801$ kcal/mol) occurs
at a similar position near the channel entrance.

Across both Cyp6p5 and Cyp6p2,
the same geographic regions segregate functionally relevant mutations
at high frequency in both species,
and the mutations themselves have convergent destabilizing effects
on protein structure.
Cytochrome P450 enzymes are central to insecticide resistance
in \textit{Anopheles},
and while most resistance at this locus has been attributed
to copy number variation,
recent evidence points to a role for nonsynonymous variation as well
\citep{Rehmann_Small_Ralph_Kern_2025}.
The convergent spatial and molecular patterns we observe at the Cyp6p locus
provide strong evidence that \textit{An. gambiae} and \textit{An. funestus}
are independently adapting to the same insecticide pressures
through parallel nonsynonymous changes in the same genes.



\begin{figure}[H]
\includegraphics[width=0.9\linewidth]{{figures/AGAP002869_mutations.pdf}}
\centering
\caption{(A) Frequency of spatially-constrained nonsynonymous Cyp6p2 mutations in \textit{An. gambiae} (left) and \textit{funestus} (right). (B) Position of spatially-constrained nonsynonymous mutations on the Cyp6p2 protein in \textit{An. gambiae} (left) and \textit{funestus} (right). Arrow indicates channel entrance.
}
\label{fig:cyp6p2}
\end{figure}

\section{Discussion}

In this study, we use gene-level population and phylogenetic approaches
to uncover the evolutionary forces acting on nonsynonymous variation
in three deeply diverged \textit{Anopheles} species.
Our results point to a scenario where diversifying selection
maintains functional nonsynonymous polymorphism within species,
and where the same genes and functional categories
that harbor this variation are also evolving under directional selection along lineages,
providing a substrate for global adaptation
when selection pressures are strong enough.


We find that nonsynonymous polymorphism is abundant in 
\textit{Anopheles} populations, and that this variation is maintained
in part by balancing selection, as indicated by the relationship 
between genic and nonsynonymous Tajima's $D$, 
and in part by local adaptation, as indicated by the relationship 
between genic and nonsynonymous $F_{ST}$.
Genes with elevated nonsynonymous polymorphism are enriched for 
sensory and metabolic functions, particularly
redox metabolism; we also identify
a number of specific genes involved in immune response and cuticle formation.
Many of these genes and functions are also evolving under 
directional selection along \textit{Anopheles} lineages,
suggesting that the similar selection pressures both maintain 
polymorphism and drive divergence in these species.
We hypothesize that much of the nonsynonymous variation we observe 
is maintained under hyperlocal adaptation, and can sweep to fixation 
under strong enough selection pressures.
Adaptation to insecticide exposure in particular appears to be 
under both diversifying and directional selection, with behavioral, 
metabolic, and cuticular genes both segregating nonsynonymous variants 
and fixing them adaptively. 
Sensory functions involved in behavioral resistance are enriched for 
nonsynonymous variation in 
\textit{An. gambiae} and \textit{minimus}, but certain genes 
involved in sensory processing have 
a significant proportion of adaptive fixations in all three lineages. 
This is also true of metabolic resistance:
 while we identify two striking examples of convergent 
local adaptation in two \textit{An. gambiae} and 
\textit{funestus} cytochrome P450s in addition 
to a number of detoxification enzymes segregating 
excess amounts of nonsynonymous polymorphism, 
overall metabolism of organophosphates is evolving
under directional selection in \textit{gambiae} and \textit{minimus}.
Genes involved in cuticle formation are also 
outliers with respect to nonsynonymous polymorphism, 
likely driven by local adaptation to desiccation and insecticide exposure, 
however cuticle formation and maintenance genes are also 
fixing adaptive variants on deeper evolutionary timescales. 
All of this points to the polymorphism maintained in 
\textit{Anopheles} populations being an abundant reservoir of 
potentially-adaptive variation relevant to the sort of
evolutionary pressures they are facing across their ranges.

The genome-wide asymptotic $\alpha$ values we estimate
($0.600$--$0.655$ across the three species)
place \textit{Anopheles} among the most adaptively evolving animal lineages studied to date.
Across a broad survey of animals,
$\alpha$ correlates positively with effective population size,
with large-$N_e$ invertebrates such as \textit{Drosophila} ($\alpha \approx 0.57$)
at the high end and small-$N_e$ vertebrates such as primates ($\alpha \approx 0.2$)
at the low end \citep{Galtier_2016}.
The \textit{Anopheles} values exceed those of \textit{D. melanogaster}
and are consistent with the large effective population sizes of these species,
confirming that selection acts efficiently to fix beneficial amino acid substitutions.
This makes the abundance of segregating nonsynonymous polymorphism all the more striking:
it persists not because selection is weak,
but despite selection being strong.
At the same time, relatively few individual genes
show significant evidence of directional selection in the per-gene MK tests,
even as the asymptotic genome-wide test indicates pervasive adaptation.
This apparent discrepancy likely reflects the strong local population structure
in \textit{Anopheles}, which generates an excess of rare nonsynonymous variants
that inflate the neutrality index at individual loci.
The asymptotic MK test accounts for this by fitting the frequency spectrum,
but individual gene-level tests lack the power to do so
and are consequently biased toward excess amino acid polymorphism.
This interpretation is consistent with the deeply negative Tajima's $D$ values
we observe across species
and with the spatial population structure evident in our $F_{ST}$ analyses.


Our phylogenetic branch model analysis provides a complementary line of evidence
connecting within-species polymorphism to between-species divergence.
By running branch models on replicate alignments of randomly sampled haplotypes,
we can distinguish genes where all sampled haplotypes show $\omega > 1$ ---
indicating that adaptive substitutions have fixed along a lineage ---
from genes where only some haplotypes do,
indicating that adaptive divergence is still segregating within the species.
This distinction is important because it reveals
that some of the nonsynonymous variation we observe within species
is not merely tolerated but is itself adaptive,
with different haplotypes at the same locus under different selective regimes.
In \textit{An. gambiae} and \textit{An. funestus},
immune genes such as LRRCT domain-containing protein
show elevated $\omega$ across all sampled haplotypes,
while others such as mannose-binding C-type lectin
show $\omega > 1$ on only a subset,
illustrating how adaptive divergence can coexist
with segregating functional polymorphism at the same locus.
These results reinforce the view that nonsynonymous variation
in \textit{Anopheles} is actively maintained by selection
rather than simply awaiting fixation or loss.

One interesting functional category that is both 
segregating high amounts of nonsynonymous variation
and evolving directionally in \textit{Anopheles} is 
genes involved in redox metabolism. 
This includes elevated nonsynonymous polymorphism in 
genes involved in glutathione metabolism, NADPH metabolism, and 
potentially the pentose phosphate pathway via enrichment 
for genes involved in purine ribonucleotide metabolism 
\citep{Kanehisa_Furumichi_Sato_Matsuura_Ishiguro-Watanabe_2024}.
\textit{Anopheles} redox metabolism is critical to 
their immune response and fecundity, 
and is also implicated in tolerance of insecticide exposure 
\citep{redox_fecundity, redox_plasmodium, redox_gutmicro}.
Given the importance of reactive oxygen species in crucial
\textit{Anopheles} biological functions and its impacts on overall
mosquito fitness, it is unsurprising that genes involved in 
redox metabolism are both maintaining nonsynonymous variation, 
possibly due to local adaptation to different pathogen and 
insecticide pressures, and also evolving directionally along 
\textit{Anopheles} lineages, with some of these variants 
becoming globally adaptive and fixing.

Beyond redox metabolism, we also uncover a number of genes and 
functional categories involved in immune response that are 
both segregating nonsynonymous variation and evolving directionally in \textit{Anopheles}.
This includes genes involved in pathogen recognition, 
such as C-type lectins and leucine-rich repeat proteins, as 
well as genes involved in \textit{Anopheles} immune signaling, 
such as CLIP-domain serine proteases.
In \textit{An. funestus}, we observe an interesting outlier 
for nonsynonymous polymorphism in a gene encoding protein TsetseEP-like
(AFUN2\_010725), a homolog of TsetseEP,
a protein involved in trypanosome defense in tsetse flies.
It is not unexpected that we observe immune genes to be both
maintaining variation and evolving directionally,
as \textit{Anopheles} are under constant pressure from a variety of pathogens
and the immune system is a common target of selection in many species.
The \textit{Anopheles} immune response is also critical to their
ability to transmit the \textit{Plasmodium} parasites,
and we observe some evidence that genes specifically involved in
\textit{Anopheles}--\textit{Plasmodium} interactions
are either maintaining variation or evolving directionally, such 
as the heparan-sulfate 6-O-sulfotransferase in \textit{An. gambiae} and
potentially Yellow-g in \textit{An. gambiae} and \textit{funestus}.
Understanding the genes required for \textit{Anopheles} vector 
competence is informative for 
potential malaria transmission-blocking strategies,
and some of these genes may be promising candidates for further 
investigation in this regard, 
particularly those that are evolving under directional selection 
as they may be more likely to be involved in critical 
\textit{Anopheles}-\textit{Plasmodium} interactions.

Human vector control efforts, particularly insecticide deployment, 
are a major selective pressure on \textit{Anopheles} populations.
Unsurprisingly, we see clear evidence of adaptation to insecticide exposure 
in all three species, with genes involved in metabolic,
behavioral, and cuticular resistance all showing signals of both 
local and directional selection.
Consistent with the relatively recent deployment of these measures 
(in evolutionary timescales) as well as the variability
in this selection pressure across space, we observe many genes relevant 
to insecticide resistance are segregating an excess of 
nonsynonymous polymorphism, particularly genes involved in sensory perception 
(behavioral resistance) and cellular responses to stress.
One example of local adaptation to potential insecticide pressure 
is identified in two cytochrome P450 genes, Cyp6p5 and Cyp6p2, 
where we observe similar spatial patterns of nonsynonymous differentiation 
in \textit{An. gambiae} and \textit{funestus} populations.
In addition to this local adaptation, we also find evidence for 
global adaptation to insecticide exposure in \textit{An. gambiae} and 
\textit{minimus}, with enrichment for genes involved in organophosphate 
metabolism amongst genes with significantly elevated
$\alpha$ values. Organophosphate insecticides are commonly used in vector 
control efforts and replaced more 
environmentally persistent organochlorine insecticides such as DDT; 
however resistance to organophosphates is becoming more 
widespread and here we observe clear evidence of adaptation to 
this class of insecticides in \textit{Anopheles} vectors 
\citep{organophosphates, organophosphate_resistance}.

One of our most striking
results is the convergent spatial differentiation in Cyp6p5 and Cyp6p2 in \textit{An. gambiae} and \textit{funestus},
with the same geographic regions segregating functionally relevant mutations in both species.
Both these genes occur at the Cyp6p locus, a cluster of cytochrome P450 genes that are associated 
with pyretheroid resistance in \textit{Anopheles} \citep{cyp6, p450}.
While much of the resistance at this locus has been attributed to copy number variation leading 
to overexpression of these detoxification enzymes, more recent evidence points to a role for nonsynonymous 
variation also being adaptive in these proteins \citep{Rehmann_Small_Ralph_Kern_2025}, and indeed 
we observe several nonsynonymous mutations in both Cyp6p5 and Cyp6p2 that are spatially constrained,
despite being at high frequency locally, and also have destabilizing effects on protein folding. 
In Cyp6p5, this includes three mutations in \textit{An. gambiae} and three in \textit{funestus} that 
occur almost exclusively in the northeast of the sampling region and are clustered at 
the N-terminus and channel entrance of the enzyme.
In Cyp6p2, we find even stronger evidence for molecular convergence, with a single mutation in either \textit{An. gambiae} 
or \textit{funestus} being strongly destabilizing and also being nearly fixed in the same geographic region.
The observation of convergent spatial and molecular patterns in these two genes provides 
an example of local adaptation to insecticide pressures, which is not entirely unsurprising given that
\textit{An. gambiae} and \textit{funestus} share a range across which they are subject to many similar environmental pressures.
Such common patterns of adaptation have been observed in other species, such as in \textit{Drosophila}, where
deeply-diverged species show parallel adaptation to environmental clines across their range \citep{Zhao_Begun_2017}.


Insecticides are of course not the only selective pressure 
acting on \textit{Anopheles} biology, and on longer evolutionary timescales we observe functional 
categories of genes evolving under directional selection that,
while presently relevant to insecticide resistance, are also 
acted on by selection pressures that predate insecticide deployment.
This includes genes involved in cuticle formation, which are 
relevant to insecticide resistance but also likely
evolving directionally in response to desiccation and other environmental 
pressures, as well as genes involved in sensory perception,
which are relevant to behavioral resistance but also critical 
for a variety of other \textit{Anopheles} 
functions such as host-seeking and mate choice.
The fact that these functional categories have been evolving under 
directional selection along \textit{Anopheles} lineages while
still segregating high amounts of nonsynonymous variation suggests that the same genes 
and pathways relevant to contemporary insecticide pressures
have been important targets of directional selection and 
local adaptation throughout \textit{Anopheles} evolutionary history,
and that the variation we observe in these genes today may 
be a future reservoir of adaptive potential for these species.

One functional category that we observe in one species 
but not others is genes involved in reproduction, which are both segregating excess 
nonsynonymous polymorphism and evolving directionally in 
\textit{An. funestus} but not \textit{gambiae} or \textit{minimus}.
One possibility is that this is due to 
cryptic species divergence within \textit{An. funestus},
 as reproductive functions are often rapidly evolving along 
 divergent lineages and may be particularly sensitive to population structure.
Indeed sympatric 'ecotypes' have been observed in \textit{An. funestus} 
populations in Burkina Faso \citep{afun_eco} and deeply-diverged 
mitochondrial lineages have been observed in sympatric \textit{An. funestus}
populations in Zambia and Tanzania \citep{afun_mito}.
Further potential evidence for cryptic species divergence may be the presence 
of large genomic inversions in \textit{An. funestus},
similar to the well-documented inversions in the \textit{An. gambiae} 
species complex that are associated with speciation
\citep{funestus, gambiae}.
While we do not find direct evidence for cryptic speciation, the 
enrichment for reproductive genes among those segregating nonsynonymous variation 
and evolving directionally in \textit{An. funestus} provides another
line of evidence that this may be the case and warrants further investigation.

More broadly, \textit{An. minimus} stands apart from the two African species
in several respects.
It shows markedly less non-neutral evolution in the McDonald-Kreitman analysis
(15.8\% of orthologs compared to $\sim$24\% in \textit{An. gambiae} and \textit{funestus}),
and its GO enrichments for nonsynonymous polymorphism center on
broad regulatory processes rather than the specific metabolic and immune pathways
identified in the African species.
Several factors may contribute to this difference.
\textit{An. minimus} occupies a distinct ecological niche
in the forested landscapes of Southeast Asia,
with different pathogen exposure and a different history of insecticide deployment
than the African vectors.
The shorter branch length between \textit{An. minimus} and \textit{An. funestus}
may also reduce power to detect lineage-specific adaptive divergence.
Additionally, our \textit{An. minimus} samples are drawn from a single country (Cambodia),
which may limit our ability to detect the spatially varying selection
that is more apparent across the continental-scale sampling of the African species.

Overall, the patterns of nonsynonymous variation and directional selection that we observe
in \textit{Anopheles} reflect both their population histories and the convergent selection
pressures they are facing across their ranges. In these highly polymorphic,
large $N_e$ species, we see evidence for both local adaptation and balancing selection
maintaining abundant nonsynonymous polymorphism that provides a reservoir of
potentially globally-adaptive variants. The same functions and processes
that segregate high amounts of nonsynonymous variation are also evolving directionally on
longer timescales, with even some of the same genes showing
parallel adaptive divergence across lineages.
This parallels the pattern observed in \textit{Drosophila}, where
species harbor comparable levels of nonsynonymous polymorphism and
also show strong evidence of adaptive protein evolution and
parallel adaptation to environmental pressures \citep{drosophila, Zhao_Wit_Svetec_Begun_2015}.
Indeed, \citet{Zhao_Begun_2017} found substantial parallelism
in \textit{Drosophila} adaptation across two timescales:
the same genes that show clinal differentiation within species
are enriched among those evolving under directional selection between species,
a finding strikingly consistent with our observation that
functional categories harboring excess nonsynonymous polymorphism in \textit{Anopheles}
are also targets of lineage-level adaptive divergence.
However, there are notable differences.
The proportion of genes evolving under directional
selection is much smaller in \textit{Anopheles} than in \textit{Drosophila} despite
much deeper divergence times,
possibly reflecting lower rates of amino acid fixation between \textit{Anopheles} lineages
and the different ecological histories of these systems.
At the same time, the high proportion of genes in \textit{Anopheles}
that harbor an excess of nonsynonymous polymorphism
is not observed in \textit{Drosophila} despite similar effective population sizes \citep{drosophila},
providing further evidence that non-neutral processes actively maintain
this variation in \textit{Anopheles}.
The specific functional targets of adaptation also differ:
where \textit{Drosophila} lineages show enrichment for
general species divergence functions such as translation and nuclear transport,
\textit{Anopheles} lineages are enriched for genes relevant to
local adaptation, immune response, and insecticide resistance.
These common patterns of selection across deeply diverged lineages highlight
the impact of parallel and convergent evolution in shaping the genomic landscape of these vectors.

Through local adaptation and balancing selection,
\textit{Anopheles} species maintain abundant nonsynonymous
polymorphism in the same functional categories that are evolving under directional selection.
Genes involved in insecticide resistance, immune response, cuticle formation, sensory perception, and redox metabolism
both segregate excess nonsynonymous polymorphism and are diverging under directional selection along \textit{Anopheles} lineages,
pointing to a scenario where standing variation provides a reservoir of potentially globally-adaptive variants.
The forces that maintain this variation and drive divergence are shared across \textit{Anopheles} lineages,
and we observe many of the same genes diverging under parallel directional selection.
Most notably, we find evidence for convergent local adaptation to insecticide pressures in \textit{An. gambiae} and \textit{funestus},
with two cytochrome P450 genes segregating molecularly similar mutations in the same geographic regions in both species.
These results demonstrate the adaptive potential of these mosquito vectors, highlighting the importance of standing
variation in their ability to respond to selection pressures across their ranges.
In \textit{Anopheles} malaria vectors, where the strongest selection pressures are often driven by human intervention,
understanding their evolutionary dynamics will be critical for anticipating and managing resistance
as vector control efforts continue to intensify.






\chapter{Conclusion}

% Understanding the evolutionary dynamics of Plasmodium and its vectors is critical for effective vector control

% spatial pop gen can help inform this

In this dissertation, I use spatial patterns of variation in  
\textit{Anopheles} malaria vectors in order to 
uncover spatial evolutionary dynamics that drive adaptation and divergence across their ranges.

In Chapter II, I present our work using correlated patterns of machine learning
prediction errors to infer common patterns of dispersal between \textit{Anopheles gambiae} 
mosquitoes and their \textit{Plasmodium falciparum} parasites. 
Being concerned about how the spatial distribution of training data impacted 
and biased our inferences, I assessed and then developed methods to reduce
these effects. When correcting for spatial bias in empirical sampling data,
correlations between \textit{Anopheles gambiae} and \textit{Plasmodium falciparum} residuals persist,
suggesting that these prediction errors reflect shared dispersal patterns 
between host and parasite.a

Given these findings, I became interested in how adaptive alleles
are spreading through these spatially-structured populations. In
Chapter III, I discover a spatial signal of positive selection through 
continuous-space simulation and use this observation to develop a method for 
detecting adaptive alleles in spatially-structured populations.
Applying this method to \textit{Anopheles gambiae}, I identify several 
novel alleles that appear to be under positive selection, including several 
at loci associated with insecticide resistance.

\input{appendices.tex}

\input{5_uothesis_bibliography}

\input{6_uothesis_end_of_document}
